\documentclass[12pt]{article}


\usepackage[top=1cm, bottom=2cm, left=2cm, right=1cm]{geometry}


\usepackage[T2A]{fontenc}
\usepackage[utf8]{inputenc}
\usepackage[english, russian]{babel}


\usepackage{amsmath, amssymb, amsthm}


\usepackage{listings}
\usepackage{xcolor}


\lstdefinestyle{haskellStyle}{
    language=Haskell,
    basicstyle=\ttfamily\small,
    keywordstyle=\color{blue}\bfseries,
    commentstyle=\color{green!50!black},
    stringstyle=\color{red},
    showstringspaces=false,
    breaklines=true,
    frame=single,
    numbers=left,
    numberstyle=\tiny\color{gray},
    stepnumber=1,
    numbersep=5pt,
    tabsize=2,
    captionpos=b
}

\lstset{style=haskellStyle}

\usepackage{graphicx}
\usepackage{hyperref}

\usepackage{siunitx}

\sisetup{
    output-decimal-marker = {,},
    group-digits = false
}


\title{О свойствах одного бесконечно вложенного радикала}
\author{В.\,Ю.\,Винник\\ \texttt{vinnik.wadim@yandex.ru}}
\date{\today}


\begin{document}

\maketitle

\begin{abstract}
Рассматривается вопрос о существовании и конечности значения бесконечно вложенного радикала с монотонно возрастающими членами. Этот вопрос сведён к сходимости
бесконечной последовательности конечных радикалов.  При этом рассмотрение значений заменяется рассмотрением функций (от величины отбрасываемого члена) и
решается вопрос об интервале, на котором эти функции ограничены сверху наперёд заданной величиной.  Также рассматриваются производные этих функий и делается
вывод о чрезвычайно малом влиянии отброшенного члена на значение выражения. Теоретическое рассмотрение подкрепляется вычислительным экспериментом.
\end{abstract}

\section{Постановка и предварительный анализ задачи}

Имеет ли определённое и конечное значение следующее выражение?
\begin{equation} \label{eq:initial}
q = \sqrt{1 + \sqrt{2 + \sqrt{3 + \sqrt{4 +\ldots}}}} .
\end{equation}

Очевидно, вопрос сводится к следующему: имеет ли предел последовательность
\begin{equation} \label{eq:val_seq}
\begin{aligned}
    q_1 & = \sqrt{1}, \\
    q_2 & = \sqrt{1 + \sqrt{2}}, \\
    q_3 & = \sqrt{1 + \sqrt{2 + \sqrt{3}}}, \\
    q_4 & = \sqrt{1 + \sqrt{2 + \sqrt{3 + \sqrt{4}}}}, \\
    \ldots
\end{aligned}
\end{equation}

\begin{equation} \label{eq:lim}
    q = \lim_{n \to \infty} {q_n} .
\end{equation}

Следует отметить, что данная последовательность определена не рекуррентно, так как значение~\(q_{n+1}\) определяется не через~\(q_n\), а отдельным выражением,
что исключает применение ряда широко известных методов.

Каждый элемент последовательности представляет собой результат отбрасывания от выражения~(\ref{eq:initial}) бесконечного хвоста, начинающегося на \(n\)-м уровне
вложенности. Поэтому представляет интерес вопрос о влиянии величины~\(x\) этого хвоста на значение всего выражения. Для последующего рассмотрения этого вопроса
следует заменить последовательность величин~\(q_n\) функциями (условимся, что всюду далее \(x \geqslant 0\))
\begin{equation} \label{eq:func_seq}
\begin{aligned}
    f_1(x) & = \sqrt{1 + x}, \\
    f_2(x) & = \sqrt{1 + \sqrt{2 + x}}, \\
    f_3(x) & = \sqrt{1 + \sqrt{2 + \sqrt{3 + x}}}, \\
    f_4(x) & = \sqrt{1 + \sqrt{2 + \sqrt{3 + \sqrt{4 + x}}}}, \\
    \ldots
\end{aligned}
\end{equation}
Тогда \(q_n = f_n(0)\) при всех натуральных~\(n\).

Примечательно, что, в отличие от последовательности величин~(\ref{eq:val_seq}), данное семейство функций допускает рекуррентное определение. При этом
оказывается естественным перед первым членом семейства добавить вырожденный 0-й, содержательно соответствующий случаю, когда отбрасываемый хвост выражения
состоит из всего выражения.
\begin{align}
    f_0(x) & = x, \label{eq:func_recur_base} \\
    f_{n+1}(x) & = f_n(\sqrt{n + x}). \label{eq:func_recur_step}
\end{align}

Вполне очевидно, что последовательность~\(q_n\) монотонно возрастает. Поэтому для доказательства существования предела~(\ref{eq:lim}) достаточно доказать
ограниченность этой последовательности сверху. Для этого, в свою очередь, выясним, на каком интервале ограничены некоторым значением функции~$f_n$.

\section{Интервалы ограниченности функций~\(f_n\)}

В качестве верхней границы возьмём число~2. Для функции~\(f_0\) тривиально: \(f_0(x) \leqslant 2\) при \(x \leqslant 2\).
Для того, чтобы функция~\(f_1\) была ограничена сверху этим значением, необходимо, чтобы подкоренное выражение не превышало~4, т.е. \(x \leqslant 3\). Для
функции~\(f_2\) имеем: внешнее подкоренное выражение не превышает~4 тогда и только тогда, когда значение под вторым корнем не превышает~9, откуда \(x \leqslant 7\).
Таким же образом легко получить, что для \(f_3(x) \leqslant 2\) необходимо и достаточно \(x \leqslant 46\). Выведем общую формулу.

Пусть для некоторого~\(n\) известно, что \(f_n(x) \leqslant 2\) т.т.т., когда \(x\leqslant r_n\). Рассмотрим функцию~\(f_{n+1}\) и найдём соответствующее
значение \(r_{n+1}\). В силу рекуррентного соотношения~(\ref{eq:func_recur_step}),
\begin{align}
f_{n+1}(x) & \leqslant 2 , \notag \\
f_n(\sqrt{n + x}) & \leqslant 2 , \notag \\
\sqrt{n + x} & \leqslant r_n , \notag \\
n + x & \leqslant r^2_n , \notag \\
r_{n+1} & = r^2_n - n . \label{eq:bound_recur}
\end{align}

Таким образом, интервал значений~\(x\), на котором функция~\(f_n\) ограничена сверху значением~2, растёт с ростом~\(n\) и, следовательно, всегда включает
\(x=0\).  Отсюда, в свою очередь, следует, что \(q_n \leqslant 2\) при всех натуральных~\(n\), что вместе с монотонностью~\(q_n\) доказывает существование
предела~(\ref{eq:lim}).

\section{К вопросу о величине отбрасываемого хвоста}

Аргумент~\(x\) функций~(\ref{eq:func_seq}) представляет собой величину хвоста, отбрасываемого от бесконечного радикала~(\ref{eq:initial}). Замена бесконечного
радикала последовательностью конечных радикалов~(\ref{eq:val_seq}) означает, как было указано выше, что отбрасываемый хвост приравнивается нулю. Между тем,
конечность значения бесконечно вложенного радикала есть именно то, что необходимо доказать. Поэтому может возникнуть сомнение в корректности приведённого выше
рассуждения.

Это сомнение несостоятельно по следующей причине: если бы вклад отбрасываемого хвоста в значение выражения~(\ref{eq:initial}) оставался существенным несмотря на
рост~\(n\), то значения~\(q_n\) отличались бы друг от друга на достаточно существенную величину, т.е. последовательность~(\ref{eq:val_seq}) не была бы
ограничена. Именно факт ограниченности~(\ref{eq:val_seq}), означает что добавление новых уровней вложенности в радикале не оказывает существенного влияния на
его значение.

Это тривиальное соображение вряд ли заслуживало бы упоминания, если бы не следующий связанный с ним аспект.  Граница~\(r_n\) интервала растёт чрезвычайно
быстро~--- как~\(2^{2^n}\) (если при достаточно больших~\(n\) пренебречь малыми членами). Следовательно, даже если бы величина~\(x\) отброшенного хвоста росла
(как функция от~\(n\)) быстро, но не быстрее, чем дважды экспоненциально, последовательность~\(f_n(x)\) всё равно имела бы предел. К этому вопросу вернёмся
позже, в разделе~\ref{sec:variations}.

Сейчас же займёмся вопросом, насколько сильно или слабо влияет величина хвоста~\(x\) на значение \(n\)-го приближения~\(f_n(x)\). Для этого следует найти
производные этих функций. Легко обнаружить, что
\begin{align*}
f_1'(x) & = \frac{1}{2\sqrt{1 + x}}, \\
f_2'(x) & = \frac{1}{2\sqrt{1 + \sqrt{2 + x}}} \cdot \frac{1}{2\sqrt{2 + x}}
= \frac{1}{4\sqrt{1 + \sqrt{2 + x}} \cdot \sqrt{2 + x}} , \\
f_3'(x) & = \frac{1}{8\sqrt{1 + \sqrt{2 + \sqrt{3 + x}}} \cdot \sqrt{2 + \sqrt{3 + x}} \cdot \sqrt{3 + x}} .
\end{align*}
Нетрудно вывести общую формулу:
\begin{equation} \label{eq:derivative}
f_n'(x) = \frac{1}{2^n \prod_{k=1}^n g_k(x)},
\end{equation}
где
\begin{equation} \label{eq:derivative_aux}
g_k(x) = \sqrt{k + \sqrt{(k + 1) + \dots + \sqrt{n + x}}}.
\end{equation}

Рассмотрим асимптотику этих производных при достаточно больших~\(n\) и~\(x\). Для начала заметим, что
\begin{align*}
g_n(x) &= \sqrt{n + x} \sim \sqrt{x}, \\
g_{n-1}(x) &= \sqrt{n-1 + g_n(x)} \sim \sqrt{\sqrt{x}} = x^{1/4}, \\
g_{n-2}(x) &= \sqrt{n-2 + g_{n-1}(x)} \sim \sqrt{x^{1/4}} = x^{1/8}, \\
&\vdots \\
g_{n-m}(x) &\sim x^{1/2^{m+1}}.
\end{align*}

В общем виде:
\[
g_k(x) \sim x^{1/2^{n-k+1}} \quad \text{при } x \to \infty.
\]

Теперь обратимся к \(p_n(x) = \prod_{k=1}^n g_k(x)\). Показатели степени \(x\) в асимптотике:
\[
\sum_{k=1}^n \frac{1}{2^{n-k+1}} = \sum_{j=1}^n \frac{1}{2^j} = 1 - \frac{1}{2^n}.
\]
Следовательно:
\[
p_n(x) \sim c_n \cdot x^{1 - 1/2^n},
\]
где \( c_n \) — некоторая константа, зависящая от \( n \).

\[
f_n'(x) = \frac{1}{2^n p_n(x)} \sim \frac{d_n}{x^{1 - 1/2^n}}, \quad d_n = \frac{1}{2^n c_n}.
\]

При фиксированном \( n \) и \( x \to \infty \):
\[
f_n'(x) \sim \frac{\text{const}}{x^{1 - \varepsilon_n}}, \quad \varepsilon_n = \frac{1}{2^n} \to 0 \text{ при } n \to \infty.
\]

В частности, при \( n \to \infty \):
\[
f_n'(x) \sim \frac{\text{const}}{x}.
\]

Таким образом, при достаточно больших~\(n\) величина~\(x\) чрезвычайно слабо (логарифмически) влияет на значение радикала~\(q_n\).

Остаётся рассмотреть вопрос в обратной постановке: как изменяется влияние фиксированного хвоста~\(x\) с ростом~\(n\). Очевидно, что оно быстро затухает, т.к.
в~\(f_n\) аргумент~\(x\) входит в степени~\(1/2^{2^n}\). Для большей строгости рассмотрим
\[
\Delta_n(x) = f_n(x) - f_n(0) .
\]
По теореме Лагранжа:
\[
\Delta_n(x) = f'_n(\xi) \cdot x, \quad \xi \in [0, x] .
\]

Производные~\(f'_n\) были ранее найдены~(\ref{eq:derivative}--\ref{eq:derivative_aux}). Очевидно, что \(g_k(x) \geqslant \sqrt{k}\) при всех \(x \geqslant 0\).
Отсюда
\[
    f'_n(x) \leqslant \frac{1}{2^n \ \prod_{k=1}^{n}\sqrt{k}} ,
\]
т.е.
\[
    f'_n(x) \leqslant \frac{1}{2^n \ \sqrt{\prod_{k=1}^{n}k}} ,
\]
или, окончательно
\[
    f'_n(x) \leqslant \frac{1}{2^n \ \sqrt{n!}} ,
\]

Это означает, что влияние~\(\Delta_n(x)\), оказываемое отброшенным хвостом~\(x\), на значение радикала~\(q_n(x)\), чрезвычайно мало: с
ростом~\(n\) оно убывает суперэкспоненциально. На практике это означает, что при приблизительном вычислении рассматриваемого предела вместо~0, как предлагалось
выше, можно было бы брать любое начальное значение~--- его вклад в итоговый результат нивелируется за небольшое число итераций.

\section{Вычислительный эксперимент}

Для проверки приведённых выше рассуждений алгоритм вычисления функции~\(f_n(x)\) был реализован на языке Haskell.
\begin{lstlisting}
nestedSqrtG :: [Double] -> Double -> Double
nestedSqrtG bs x = foldl h x (reverse bs) where
    h u v = sqrt (u + v)

nestedSqrtF :: Int -> Double -> Double
nestedSqrtF n x = nestedSqrtG (map fromIntegral [1..n]) x

nestedSqrt n = nestedSqrtF n 0
\end{lstlisting}

Здесь функция \lstinline{nestedSqrtG} вычисляет значение вложенного радикала, если известна последовательность чисел~\(b_k\), стоящих первыми слагаемыми под каждым
корнем, и при известном значении~\(x\), которое подставляется вторым слагаемым в последний (наиболее внутренний) корень. Она очевидно представляет собой свёртку
последовательности~\(b\) по операции \(h(u, v) = \sqrt{u + v}\) с начальным значением~\(x\).
Функция \lstinline{nestedSqrtF} подставляет вместо~\(b\) последовательность чисел от~1 до~\(n\).
Наконец, функция \lstinline{nestedSqrt} фиксирует начальное значение~\(x\) равным~0.

Убедимся сначала в быстрой сходимости~\(f_n(x)\) с ростом~\(n\) при различных значениях \(x \in \{0, 10, 100\}\), см. табл.~\ref{tab:converge}. Из таблицы
видно, что значительное различие начальных значений уже при десятом приближении проявляется лишь в шестом знаке после запятой. При нулевом начальном значении,
что соответствует последовательности~(\ref{eq:val_seq}), наблюдается быстрая стабилизация десятичных знаков одного за другим.

\begin{table}[htbp]
\centering
\caption{Сходимость \(f_n(x)\)}
\label{tab:converge}
\begin{tabular}{|l|S[table-format=2.16]|S[table-format=2.16]|S[table-format=2.16]|}
\hline
{\(n\)} & {\(f_n(0)\)} & {\(f_n(10)\)} & {\(f_n(100)\)} \\
\hline
1     &  1,0                  &  3,3166247903554      &  10,04987562112089  \\
2     &  1,5537739740300374   &  2,112842070562245    &  3,3315919525599287 \\
3     &  1,7122650649295326   &  1,8351034149944543   &  2,117906043670099  \\
4     &  1,7487627132551438   &  1,7731153394806531   &  1,8366750797310631 \\
5     &  1,7562384875823431   &  1,7606505920882154   &  1,7735307361095183 \\
6     &  1,7576412350415822   &  1,758382204066681    &  1,760745403334578  \\
7     &  1,7578856460964372   &  1,7580022858671953   &  1,7584014463169158 \\
8     &  1,7579255575682606   &  1,7579429116634482   &  1,7580058397112581 \\
9     &  1,7579317105145658   &  1,757934166521764    &  1,757943518702261  \\
10    &  1,757932611393831    &  1,7579329436783082   &  1,7579342635288981 \\
\hline
\end{tabular}
\end{table}

Рассмотрим выражение языка Haskell
\begin{lstlisting}
let qs = map nestedSqrt [1..] in
    length (takeWhile (/=0.0) (zipWith (-) (tail qs) qs))
\end{lstlisting}
Оно строит последовательность значений функции \lstinline{nestedSqrt} на всех натуральных аргументах от~1 до бесконечности, т.е. последовательность
приближений~(\ref{eq:val_seq}), затем находит разности пар соседних значений и находит номер приближения, которое отличается от следующего на величину, меньшую,
чем поддерживаемая вычислительной машиной точность вещественных чисел (порядка~\(10^{-16}\)). Вычисление этого выражения даёт значение~19, т.е. 19 уровней
вложенности хватает, чтобы вычислить искомый предел с максимально возможной для машины точностью.

Подставив в предыдущее выражение функцию \lstinline{nestedSqrtF} с различными начальными значениями~\(x\), можно обнаружить, что при \(x=100\) для достижения максимально
возможной точности требуется лишь на одно приближение больше, а при \(x=10000\)~--- на два приближения. Этим подтверждаются сделанные выше выводы о чрезвычайно
быстрой сходимости~\(f_n(x)\) при любых~\(x\) и о чрезвычайно слабом влиянии начального значения~\(x\) на значение~\(f_n(x)\) при достаточно больших~\(n\).

\section{Обобщения и постановка новых задач}\label{sec:variations}

Рассмотрим данную задачу в слегка изменённой формулировке. Пусть дана последовательность чисел~\(b_n\), и пусть
\begin{equation} \label{eq:var_initial}
q = \sqrt{b_1 + \sqrt{b_2 + \sqrt{b_3 + \sqrt{b_4 +\ldots}}}} .
\end{equation}
Теперь можно для различных последовательностей~\(b\) решать вопрос о существовании этого значения, а также ставить вопрос о том, при каких условиях, наложенных
на последовательность~\(b\), это значение существует.

Ясно, что исходной задаче~(\ref{eq:var_initial}) соответствует \(b_n = n\).
Ещё один хорошо исследованный частный случай получим, взяв \(b_n = 1\) при всех~\(n\). Легко видеть, что в этом случае искомое значение~\(q\) удовлетворяет
равенству
\[
    q = \sqrt{q + 1} ,
\]
откуда получаем
\[
    q^2 - q - 1 = 0 ,
\]
т.е. \(q\) есть золотое сечение.

Выше была выведена формула~(\ref{eq:bound_recur}) для границы интервала, на котором функции~\(f_n(x)\) ограничены сверху значением~2. Представляет интерес
обобщение этой формулы на случай произвольной последовательности~\(b\).
Подобным же образом следует найти, для каких последовательностей~\(b\) остаётся в силе вывод о чрезвычайно слабой зависимости~\(f_n(x)\) от~\(x\).
Исходя из дважды экспоненциального разрастания интервала~\(r\) и сверхэкспоненциального затухания хвоста в исходной задаче можно ожидать, что
выражение~(\ref{eq:var_initial}) будет иметь значение, если последовательность~\(b_n\) растёт медленнее, чем~\(O(2^{2^n})\), однако этот вопрос требует
отдельного рассмотрения.

\end{document}
